\documentclass[11pt]{article}
\usepackage[margin=1in]{geometry}
\usepackage{amsmath,amssymb}
\usepackage{booktabs}
\usepackage{longtable}
\usepackage{hyperref}

\newcommand{\lamV}{\lambda_V}
\newcommand{\lamS}{\lambda_S}
\newcommand{\lamhat}{\hat\lambda}

\title{Asymmetric Capital Buffer Control:\\ Empirical Companion}
\author{}
\date{11 August 2026}

\begin{document}
\maketitle

\begin{abstract}
\noindent
This note is the empirical companion to \emph{Asymmetric Capital Buffer
Control} (PROOFS\_v2); numbered results cited here --- Theorem~1,
Proposition~11, Remark~1 --- refer to that document. It
states what the data say about the variance-asymmetry parameter $\lamV$,
against the null of a risk-neutral bank facing the same regulation. Every
figure below is regenerated from the scripts named in
Section~\ref{sec:repro}; none is carried from an earlier draft. The
primary analysis is the pre-registered design, reported exactly as run,
including its attrition; a secondary analysis with an alternative
reference construction is specified in Section~\ref{sec:secondary} with
results deliberately pending its own validation.
\end{abstract}

\section{What is being estimated, and how}
\label{sec:estimand}

Theorem~1 makes $\lamV^4$ the ratio of conditional variances of the log
capital-buffer state deep in deficit to deep in surplus. Two facts about
estimating it, both established by simulation
(\texttt{verify\_egarch\_recovery.py}), fix the design.

\emph{First, $\theta$ and $\lamV$ are not jointly identified from
routine-sized excursions.} With abundant data ($T=5000$) and the state
ranging over $[-1.9, 1.4]$ against $\theta=1$, the profile log-likelihood
in $\theta$ is flat to within about one point from $\theta=0.7$ to
$\theta=5$ while the implied $\lamV$ runs from $1.37$ to $6.17$: many
$(c,\theta)$ pairs are close to observationally equivalent, so $\theta$
must come from outside the estimation. The practical consequence adopted
here is an estimator that does not use $\theta$ at all: $\lamhat$ is the
fourth root of the ratio of within-regime innovation variances, each from
a within-regime AR(1) residual, so only the regime classification and the
variance ratio enter.

\emph{Second, conditional on that choice, recovery is clean.} At the null
(true $\lambda=1$, $T=300$, twenty independent replications) the mean
bias is $+0.085$ with standard deviation $0.119$, and the bias stays
small and one-sided across a $(\phi,T)$ grid from $\phi=0.5$ to $0.95$
and $T=150$ to $800$ --- shrinking to $+0.018$ at $(\phi,T)=(0.95,800)$.
There is no free trend parameter whose mis-estimation could manufacture
asymmetry, which is the failure mode that produced $\lamhat\approx 2$ at
a true null in a sibling project's estimator. The small positive null
bias means the estimator, if anything, overstates asymmetry slightly ---
so it works against, not for, the direction of the finding reported
below.

\section{The null: what a risk-neutral bank would already show}
\label{sec:null}

The comparison value is not $\lamV=1$. A risk-neutral bank
($\lamS=1$) facing the regulatory cap already generates asymmetric
conditional variance, because the cap tightens exposure as the buffer
erodes; the honest null is therefore $\lamV$ computed from the solved
model itself. Two routes exist and only one survives.

The \emph{definitional} route --- Theorem~1's ratio evaluated at the
saturation limits --- is a constant of cap geometry,
$\lamV = (\kappa_1/\kappa_3)^{1/2} = 2.599$ at the baseline, with no
dependence on $\lamS$ or any preference or cost parameter
(Proposition~11). It cannot serve as a benchmark for data that never
reach the saturation limits, and it carries no information about
$\lamS$.

The \emph{estimator} route applies the same $\theta$-free estimator to
paths simulated from the solved model. Computed on the
convergence-certified solves (Remark~1) at
$\lamS\in\{1.00, 1.10, 1.25, 1.50, 1.75, 2.00\}$, under three
conventions:

\begin{center}
\begin{tabular}{lcccccc}
\toprule
& \multicolumn{6}{c}{$\lamS$} \\
\cmidrule(lr){2-7}
Convention & 1.00 & 1.10 & 1.25 & 1.50 & 1.75 & 2.00 \\
\midrule
(a) RWA$\sim$L, split at stationary mean & 1.3320 & 1.3269 & 1.2892 & 1.2724 & 1.2477 & 1.2347 \\
(b) RWA$\sim$A, split at stationary mean & 1.4762 & 1.4764 & 1.4259 & 1.4038 & 1.3723 & 1.3545 \\
(a$'$) RWA$\sim$L, split fixed at $R$ & 1.3100 & 1.3423 & 1.3590 & 1.3721 & 1.3967 & 1.4056 \\
\bottomrule
\end{tabular}
\end{center}

Three properties of this table matter. The $\lamS=1$ column is the
risk-neutral null under each convention: $1.31$--$1.48$. The
\emph{direction} in which $\lamV$ moves with $\lamS$ is
convention-dependent: both moving-threshold conventions decrease
(slopes $-0.10$ and $-0.13$ per unit $\lamS$), while the fixed-threshold
convention increases ($+0.09$); the likely mechanism is that under a
moving threshold the split point itself shifts with $\lamS$ (from
$1.30$ to $1.43$ across the grid), confounding the $\lamS$ effect with a
changing comparison point. The panel estimator of
Section~\ref{sec:primary} classifies regimes against a moving,
history-based threshold, so rows (a)--(b) are the relevant comparison
for it, and under them $\lamV$ \emph{falls} as asymmetry rises. Finally,
the smallest value anywhere in the table is $1.2347$: no admissible
$\lamS$ in the solved range produces $\lamV$ below that, under any
convention.

\section{Primary analysis: the pre-registered design, as run}
\label{sec:primary}

Run on 11~August~2026 against the live IMF FSI panel under the stated
universe rule (\texttt{panel\_universe.py}: every sovereign or
dependent-territory area the source reports, no history pre-screen;
152 areas returned, 151 sovereign, one regional aggregate excluded by
name). Of the 151, 114 carry enough history to attempt estimation and 37
are too short. The estimator produces a usable $\hat\lambda$ for
\textbf{33 countries}: median $\mathbf{1.038}$, mean $1.132$, interquartile
range $[0.963,\,1.288]$, fraction above one $0.61$. Ten of the 33 sit at
or above the lowest model value; thirteen sit below one. One estimate
(Equatorial Guinea) carries a within-regime AR coefficient of $-3.9$ and
should be treated as unreliable regardless of its regime counts.

The headline is robust to the universe correction: the previous,
unexplained 77-country list gave median $1.0395$ on $n=26$; the full
sovereign universe, roughly twice the size, gives $1.0379$ on $n=33$.
Doubling the universe moved the median by $0.002$.

\subsection{The attrition, stated}

Of the 114 attempted, 81 do not yield an estimate, and the two mechanisms
must be reported separately because only one is curable by threshold
choice. \textbf{Structural} exclusions --- zero quarters in one regime,
which no per-regime minimum can repair --- account for 24 countries, of
which eighteen are advanced or EU economies: France (39 surplus quarters,
0 deficit), the United Kingdom (44/0), Belgium (49/0), Sweden (39/0),
Austria (51/0), Ireland (48/0), Israel (64/0), the Netherlands (56/0),
Norway (52/0), and the remainder in central and northern Europe.
\textbf{Near-miss} exclusions --- a nonzero smaller regime short of the
per-regime minimum of 15 --- include Switzerland, Spain and Denmark at 2
deficit quarters each, Luxembourg at 1, Greece at 3, Portugal at 6, and
the United States at 14.

The sharpened form of the selection finding is therefore: \emph{under a
one-sided expanding-mean reference, the deficit regime is unobservable in
every advanced banking system in the post-Basel~III sample.} Not thinly
observed --- unobservable. A capital ratio that trends upward from the
start of its reported history never falls below its own expanding mean,
and the advanced economies are precisely the trending ones. An
advanced-economy-only panel under this construction would have $n=4$
(Canada, Estonia, Macao, San Marino). The measurable population is
selected on the dynamics of the variable being measured, and the
population of greatest interest for the theory is the one selected out.
Section~\ref{sec:secondary} exists for exactly this reason.

\subsection{Reading the result}

The comparison must be made on the measured scale. The estimator is
biased toward one, and the bias does not vanish with sample size: a
process with true $\lambda=1.30$ measures $1.16$ at $T=60$ quarters and
still $1.14$ at $T=20{,}000$, because observations near the
classification boundary are assigned to a regime while their innovations
reflect a mixture of both. It is a different estimand, not a small-sample
artefact. The model's population value (Section~2) is therefore not the
right benchmark for the panel median; the right benchmark is what the
panel's own estimator returns when applied to the model
(\texttt{results/lambda\_RN\_measured.py}: the solved model simulated at
quarterly frequency over 60-quarter windows, estimated by the identical
code path, 400 paths per parameter value).

On that scale, the model class spans median measured values of
$\mathbf{1.334}$ (risk-neutral, convention (a)) down to $\mathbf{1.152}$
(the most asymmetric admissible parameter, $\lambda_S=2$), with
convention (b) slightly higher throughout, and a per-path sampling
standard deviation of $0.16$--$0.20$.

Against this the panel median of $1.038$ sits below every cell of the
grid: $0.30$ below the risk-neutral benchmark (roughly six to seven
standard errors of a median of 33) and $0.11$ below even the most
asymmetric admissible model (roughly two to three standard errors). The
direction is unambiguous under both conventions and every admissible
$\lambda_S$: \emph{measured deficit-state volatility in the usable panel
is lower than a preference-free regulated bank would be measured at}.
The magnitude, however, is modest once stated on the measured scale, and
no claim is made here that a preference parameter is identified from it:
13.2\% of individual simulated model banks measure at or below $1.038$
over sixty quarters, so any single country at that level is unremarkable
and the evidence rests entirely on the median of 33. The population
value of the gap is roughly $0.20$ if the attenuation mapping is
inverted, but that inversion is model-dependent and is not used for any
claim.

Two further cautions. First, the measured-scale model values are far
flatter in $\lambda_S$ ($1.334$ to $1.152$ across the whole admissible
range) than the population values, so the data compress exactly the
variation that would be needed to read $\lambda_S$ back from a variance
ratio --- consistent with the identification analysis in Section~2, and
one more reason the committed thresholds, not the variance ratio, are
the identifying moments if identification is ever attempted. Second, the
usable panel is the mean-reverting, largely emerging-market subsample
that the reference construction itself selected; the undershoot is a
fact about that population, generalised to no other.

\section{Secondary analysis: alternative reference (specified;
results pending validation)}
\label{sec:secondary}

The attrition of Section~\ref{sec:attrition} motivates a second
reference construction, recorded here \emph{before} its results are
computed, for four reasons: the theory's reference is a threshold, not a
country's own trend, so ``below your own past'' is not the estimand
Theorem~1 defines; the surviving sample is selected on behaviour
plausibly correlated with the quantity measured; the excluded countries
are exactly the deep-banking-system economies the theory most concerns;
and the model null of Section~\ref{sec:null} is computed from a
stationary distribution, so the comparison currently holds only on the
subsample where stationarity approximately obtains.

\emph{Construction.} Per country, the reference is the full-sample mean
of the log ratio, $R_i = \mathrm{mean}(\log\text{ratio}_{i,\cdot})$,
constant over time; classification and the $\theta$-free estimator are
otherwise unchanged. A trending series lies below its full-sample mean
early and above it late, so both regimes exist for every country and the
13 zero-deficit countries re-enter.

\emph{Known costs, stated in advance.} The construction uses look-ahead
--- the reference at date $t$ involves data after $t$ --- and for a
trending series the regime label partially proxies for calendar time, so
$\lamhat$ may pick up secular changes in volatility rather than
state-dependence. Both effects are testable in simulation, and the
design is therefore gated exactly as the primary was: it runs on
simulated paths first --- recovery against the true-$z$ benchmark, a
clean null with no manufactured asymmetry, and additionally a
\emph{trended} null (upward-drifting ratio, symmetric shocks), which is
the case the look-ahead concern targets --- and touches real data only
if it passes all three. No bandwidth or window parameter is introduced,
so there is no tuning margin between this specification and its result.

\emph{Interpretation rule, fixed now.} If the secondary agrees with the
primary (median well below the null range), the finding generalises
beyond the mean-reverting subsample. If it disagrees, the primary result
is evidence about the reference construction, not about banks, and
neither number should be quoted without the other.

\section{What this note does not establish}
\label{sec:limits}

The calibration behind the null is the project's stylised parameter set,
not fitted to bank data. The pre-committed $\phi^-$ test, with its fixed
treatment of the nine ROE${}<g$ countries, remains unexecuted, as does
the bank-level feasibility check
(\texttt{empirical\_design\_exclusions}). The RWA$\sim$A versus
RWA$\sim$L choice is settled by the definitional argument that Basel's
risk weights are asset-side, but has not been checked against the IMF's
construction of FSKRC\_PT; the primary reading does not depend on it,
since the median sits below the null under both. And whether the solved
model's constructed value function is the value function --- the
verification theorem --- remains open (PROOFS\_v2, ``What is not
established''), so the null values here are properties of the certified
numerical solution.

\section{Reproduction}
\label{sec:repro}

Panel: \texttt{03\_empirical/design/expanding\_mean\_panel.py} (live IMF
FSI fetch; requires network and \texttt{statsmodels}); run
11~August~2026, output reproduced in Section~\ref{sec:primary}. Null
curve: \texttt{02\_numerical/run\_lambda\_sweep.sh} at
\texttt{max\_iter} $\ge 2000$, then
\texttt{run\_lambda\_RN\_curve.sh} against
\texttt{lambda\_RN\_panelconv.m}; the certified solves and the resulting
\texttt{lambda\_V\_curve.csv} are in \texttt{03\_empirical/results/}.
Estimator properties: \texttt{01\_theory/verify\_egarch\_recovery.py}
(Parts A, B, E quoted in Section~\ref{sec:estimand}). The secondary
analysis of Section~\ref{sec:secondary} has, deliberately, no results
file yet; its validation harness extends the same script.

\end{document}
